\documentclass[11pt]{article}

% ---------- Packages ----------
\usepackage[T1]{fontenc}
\usepackage[margin=0.85in]{geometry}
\usepackage{booktabs}
\usepackage{array}
\usepackage{longtable}
\usepackage{enumitem}
\usepackage{xcolor}
\usepackage{titlesec}
\usepackage{hyperref}
\usepackage{fancyhdr}
\usepackage{parskip}
\usepackage{ragged2e}

% ---------- Colors ----------
\definecolor{jhublue}{HTML}{002D72}
\definecolor{jhugold}{HTML}{68ACE5}
\definecolor{diffLow}{HTML}{2E7D32}
\definecolor{diffMed}{HTML}{E65100}
\definecolor{diffHigh}{HTML}{B71C1C}

% ---------- Section styling ----------
\titleformat{\section}{\Large\bfseries\color{jhublue}}{\thesection}{0.8em}{}
\titleformat{\subsection}{\large\bfseries\color{jhublue!85!black}}{\thesubsection}{0.8em}{}
\titlespacing*{\section}{0pt}{1.4em}{0.8em}
\titlespacing*{\subsection}{0pt}{1.1em}{0.5em}

% ---------- Header / Footer ----------
\pagestyle{fancy}
\fancyhf{}
\renewcommand{\headrulewidth}{0.4pt}
\fancyhead[L]{\small EN.553.439/639 --- Time Series Analysis}
\fancyhead[R]{\small Spring 2026 Course Breakdown}
\fancyfoot[C]{\small\thepage}

% ---------- Difficulty macros ----------
\newcommand{\Low}{\textcolor{diffLow}{\textbf{Low}}}
\newcommand{\Med}{\textcolor{diffMed}{\textbf{Medium}}}
\newcommand{\High}{\textcolor{diffHigh}{\textbf{High}}}
\newcommand{\MedHigh}{\textcolor{diffMed}{\textbf{Med}}--\textcolor{diffHigh}{\textbf{High}}}

\newcommand{\module}[2]{\subsection*{Module #1: #2}}

\hypersetup{
  colorlinks=true,
  linkcolor=jhublue,
  urlcolor=jhublue,
  citecolor=jhublue,
  pdftitle={EN.553.439/639 Time Series Analysis -- Course Breakdown},
  pdfauthor={Course Breakdown}
}

\setlist[itemize]{leftmargin=1.4em, itemsep=0.15em, topsep=0.2em}
\setlist[enumerate]{leftmargin=1.6em, itemsep=0.15em, topsep=0.2em}

\begin{document}

% ================= TITLE =================
\begin{center}
  {\Huge\bfseries\color{jhublue} Time Series Analysis}\\[0.3em]
  {\Large EN.553.439 / EN.553.639}\\[0.2em]
  {\large Full Course Breakdown --- Modules, Learning Objectives \& Pacing}\\[0.4em]
  {\normalsize Spring 2026 \(\cdot\) Johns Hopkins University \(\cdot\) Dr.\ Sergey Kushnarev}\\[0.1em]
  {\small Derived from the course syllabus (Jan.\ 20, 2026 draft) and \emph{Cryer \& Chan, Time Series
  Analysis with Applications in R}, 2nd ed.}
\end{center}

\vspace{0.6em}
\hrule
\vspace{1em}

% ================= 1. OVERVIEW =================
\section{Course Overview}

Time Series Analysis (EN.553.439/639) covers time series methods from both the \textbf{time-domain}
and \textbf{frequency-domain} perspectives. The course moves from descriptive statistics and trend
estimation through the full Box--Jenkins ARIMA modeling cycle (identification, estimation,
diagnostics, forecasting), then extends to seasonal, regression-based, and heteroskedastic
(ARCH/GARCH) models. It covers approximately the first ten chapters of Cryer \& Chan, plus selected
material from Chapter 12.

\begin{itemize}
  \item \textbf{Credits:} 3 (EQ) \qquad \textbf{Format:} Two 75-minute lectures/week (MW, 4:30--5:45 PM)
  \item \textbf{Grading:} 6 homework assignments (25\% total) + 3 non-cumulative in-class midterms
        (25\% each = 75\%); final grades are curved
  \item \textbf{Assessment philosophy:} Each midterm covers a distinct third of the material
        (non-cumulative), so the course naturally decomposes into \textbf{three instructional units}
        (see Section~4)
\end{itemize}

% ================= 2. LEARNING OUTCOMES =================
\section{Course-Level Learning Outcomes}

By the end of the course, a successful student will be able to:

\begin{enumerate}
  \item Characterize a time series using mean, autocovariance, and autocorrelation functions, and
        determine whether a stochastic process is (weakly/strictly) stationary.
  \item Formulate deterministic- and stochastic-trend models and apply smoothing/detrending techniques.
  \item Represent AR, MA, ARMA, and ARIMA processes using backshift-operator notation and derive their
        theoretical ACF/PACF.
  \item Identify a tentative model order for a real data set using sample ACF, PACF, and EACF.
  \item Estimate model parameters via Yule--Walker, (conditional) least squares, and maximum likelihood,
        and assess estimator adequacy.
  \item Diagnose a fitted model via residual analysis and formal goodness-of-fit tests, and compare
        competing models.
  \item Construct point forecasts and prediction intervals for stationary and nonstationary series, and
        update forecasts as new data arrive.
  \item Build multiplicative SARIMA models for seasonal series and carry them through the full modeling
        cycle.
  \item Apply regression-based techniques --- intervention analysis, outlier detection, pre-whitening
        --- to series with exogenous effects, and explain the danger of spurious correlation.
  \item Recognize conditional heteroskedasticity (volatility clustering) and specify basic ARCH/GARCH
        models.
  \item Interpret a time series in the frequency domain via the periodogram and spectral density.
\end{enumerate}

% ================= 3. PREREQUISITES =================
\section{Prerequisite Knowledge Map}

Formal prerequisite: \textbf{(553.310+ OR 553.311+ OR 553.420+ OR 553.421+ OR 553.620+) AND
(AS.110.201+ OR AS.110.212+ OR 553.291+)}. In practice this means one probability/mathematical-statistics
course \emph{and} one linear algebra course. The table below translates this into the specific skills
the course draws on from day one.

\renewcommand{\arraystretch}{1.25}
\begin{longtable}{p{3.6cm} p{3.6cm} p{7.1cm}}
\toprule
\textbf{Prerequisite Area} & \textbf{Typical Source} & \textbf{Specific Skills Needed in TSA} \\
\midrule
\endhead
Probability \& mathematical statistics &
  553.310/311/420/421/620 (probability sequence) &
  Random variables, expectation/variance/covariance, conditional expectation, basic stochastic-process
  ideas (independence, i.i.d.\ sequences); needed immediately for the ACVF/ACF and for defining
  stationarity. \\
Linear algebra &
  AS.110.201/212 or 553.291 &
  Matrix/vector operations, eigenvalues and eigenvectors, quadratic forms, systems of linear equations;
  needed to solve Yule--Walker equations, understand spectral decompositions, and (optionally)
  state-space formulations. \\
Calculus \& optimization &
  Multivariable calculus (implicit) &
  Differentiation, chain rule, multivariable optimization; needed to derive least-squares and maximum
  likelihood estimators for ARMA parameters. \\
Regression \& hypothesis testing &
  Introductory applied statistics (implicit) &
  Ordinary least squares, residual analysis, confidence/hypothesis tests; needed for trend estimation,
  model diagnostics, and the regression-based methods module. \\
Basic R programming &
  Not formally required, but expected &
  Reading in data, fitting models with built-in functions, basic plotting; the required text uses R
  throughout and homework will assume comfort with it. \\
\bottomrule
\end{longtable}

% ================= 4. STRUCTURE =================
\section{Course Structure: Three-Unit Design}

The three non-cumulative midterms partition the semester into three roughly equal instructional
units. Dates below are the instructor-specified exam dates, cross-checked against the published JHU
Spring 2026 academic calendar (classes begin Tue., Jan.\ 20; Spring Break is Mon.--Fri., Mar.\ 16--20;
last day of classes is Mon., Apr.\ 27).

\renewcommand{\arraystretch}{1.3}
\begin{longtable}{p{1.6cm} p{3.0cm} p{6.3cm} p{2.9cm}}
\toprule
\textbf{Unit} & \textbf{Weeks} & \textbf{Focus} & \textbf{Milestone} \\
\midrule
\endhead
I &
  Weeks 1--5 (Jan.\ 21 -- Feb.\ 18) &
  Foundations: descriptive tools, stationarity, deterministic/stochastic trends &
  \textbf{Midterm 1} \newline Wed., Feb.\ 18 \\
II &
  Weeks 6--9 (Feb.\ 23 -- Mar.\ 25, spans Spring Break Mar.\ 16--20) &
  Linear stochastic models (ARMA/ARIMA) and the model-identification \& estimation toolkit &
  \textbf{Midterm 2} \newline Wed., Mar.\ 25 \\
III &
  Weeks 10--13 (Mar.\ 30 -- Apr.\ 22) &
  Diagnostics, forecasting, seasonal (SARIMA) models, regression-based methods &
  \textbf{Midterm 3} \newline Wed., Apr.\ 22 \\
Wrap-up &
  Week 14 (Apr.\ 27, last day of classes) &
  Heteroskedastic models (ARCH/GARCH) and frequency-domain/nonlinear topics ``as time allows'' &
  Review \\
\bottomrule
\end{longtable}

\textit{Note: the university calendar's own week numbering skips the Spring Break week rather than
counting it, which is why Weeks 6--9 span a five-calendar-week range.}

% ================= 5. MODULE BREAKDOWN =================
\section{Module-by-Module Breakdown}

The syllabus's 13 numbered topics map to modules below. ``Ch.'' gives the approximate corresponding
chapter in Cryer \& Chan; treat these as approximate since the syllabus states the course covers
material ``approximately'' following the first ten chapters plus parts of Chapter 12.

\subsection*{Quick-Reference Table}
\renewcommand{\arraystretch}{1.3}
\begin{longtable}{p{0.5cm} p{4.3cm} p{1.0cm} p{6.2cm} p{1.6cm}}
\toprule
\textbf{\#} & \textbf{Module} & \textbf{Ch.} & \textbf{Key Focus} & \textbf{Difficulty} \\
\midrule
\endhead
1 & Fundamental Tools & 1--2 & Mean, variance, covariance, ACVF, ACF (population \& sample) & \Low \\
2 & Stationarity & 2 & Strict vs.\ weak stationarity, white noise & \Med \\
3 & Trends & 3 & Deterministic vs.\ stochastic trends, smoothing, detrending & \Med \\
4 & Linear (ARMA) Models & 4 & General linear process, AR/MA/ARMA, backshift operator, stationarity/invertibility & \High \\
5 & Nonstationary Models & 5 & Differencing, ARIMA, variance-stabilizing transforms & \Med \\
6 & Model Specification & 6 & ACF/PACF/EACF-based identification & \High \\
7 & Parameter Estimation & 7 & Yule--Walker, (conditional) LS, MLE & \High \\
8 & Model Diagnostics & 8 & Residual analysis, portmanteau tests, model comparison & \Med \\
9 & Forecasting & 9 & Minimum-MSE forecasts, prediction intervals, updating & \MedHigh \\
10 & Seasonal (SARIMA) Models & 10 & Seasonal differencing, multiplicative SARIMA & \High \\
11 & Regression-Based Methods & 11 & Intervention analysis, outlier detection, spurious correlation, pre-whitening & \MedHigh \\
12 & Heteroskedastic Models & 12 (part) & ARCH, GARCH, volatility clustering & \High \\
13 & Additional/Advanced Topics & 12 (part) & Nonlinear TS, spectral/frequency-domain methods & \High \\
\bottomrule
\end{longtable}

\subsection*{Detailed Module Descriptions}

\module{1}{Fundamental Tools}
\textbf{Concepts:} time series as a realization of a stochastic process; mean function; variance and
covariance; autocovariance function (ACVF); autocorrelation function (ACF); sample vs.\ population
versions.\\
\textbf{Learning objectives --- students will be able to:}
\begin{itemize}
  \item Define a time series as a realization of an underlying stochastic process.
  \item Compute the sample mean and sample ACF of a data set and read a correlogram.
  \item Distinguish population quantities from their sample estimates.
\end{itemize}
\textbf{Difficulty:} \Low\ --- notation-heavy but conceptually approachable given the probability
prerequisite.

\module{2}{Stationarity}
\textbf{Concepts:} strict stationarity; weak (second-order/covariance) stationarity; properties of the
ACVF under stationarity; white noise.\\
\textbf{Learning objectives:}
\begin{itemize}
  \item Distinguish strict from weak stationarity and identify which is required for a given result.
  \item Verify whether a simple process (e.g., a moving average of white noise) is stationary.
  \item Explain why stationarity is a prerequisite for most of the course's modeling machinery.
\end{itemize}
\textbf{Difficulty:} \Med\ --- the strict-vs-weak distinction is conceptually subtle and a common
early stumbling block.

\module{3}{Trends}
\textbf{Concepts:} deterministic vs.\ stochastic trends; least-squares trend estimation; smoothing
methods (moving average, exponential smoothing); detrending.\\
\textbf{Learning objectives:}
\begin{itemize}
  \item Distinguish deterministic-trend from stochastic-trend (random-walk-type) behavior.
  \item Fit and evaluate a deterministic trend model via least squares.
  \item Apply a smoothing method to separate trend from noise.
\end{itemize}
\textbf{Difficulty:} \Med.

\module{4}{Linear (ARMA) Models}
\textbf{Concepts:} general linear process (Wold-type representation); AR($p$), MA($q$), ARMA($p,q$);
the backshift operator; stationarity and invertibility conditions via characteristic roots.\\
\textbf{Learning objectives:}
\begin{itemize}
  \item Derive the theoretical ACF/PACF of AR(1), AR(2), and MA(1) processes.
  \item Determine stationarity and invertibility using the roots of the characteristic equation.
  \item Write AR/MA/ARMA models in backshift-operator notation and manipulate that notation.
\end{itemize}
\textbf{Difficulty:} \High\ --- typically the first major conceptual hurdle; root-finding and operator
algebra are unfamiliar to many students at this point.

\module{5}{Nonstationary Models}
\textbf{Concepts:} differencing; random walk; ARIMA($p,d,q$); variance-stabilizing transformations
(log, Box--Cox).\\
\textbf{Learning objectives:}
\begin{itemize}
  \item Determine the order of differencing $d$ needed to induce stationarity.
  \item Fit an ARIMA model to a nonstationary series.
  \item Choose and apply an appropriate variance-stabilizing transformation.
\end{itemize}
\textbf{Difficulty:} \Med.

\module{6}{Model Specification / Identification}
\textbf{Concepts:} sample ACF and PACF as identification tools; the Extended ACF (EACF); tentative
order selection.\\
\textbf{Learning objectives:}
\begin{itemize}
  \item Read sample ACF/PACF plots to propose candidate AR, MA, or ARMA orders.
  \item Use the EACF table when ACF/PACF patterns are ambiguous.
  \item Justify a model choice by combining multiple identification tools.
\end{itemize}
\textbf{Difficulty:} \High\ --- this is the most judgment-driven part of the course (``art vs.\
science''), and students often find the lack of a single deterministic rule uncomfortable.

\module{7}{Parameter Estimation}
\textbf{Concepts:} method-of-moments/Yule--Walker estimation; least squares and conditional least
squares; maximum likelihood estimation (MLE); basic estimator properties.\\
\textbf{Learning objectives:}
\begin{itemize}
  \item Set up and solve the Yule--Walker equations for AR parameter estimates.
  \item Distinguish conditional from unconditional (exact) least squares/likelihood.
  \item Explain, at a working level, how the ARMA likelihood is constructed and optimized.
\end{itemize}
\textbf{Difficulty:} \High\ --- requires calculus-based optimization and careful bookkeeping of which
likelihood approximation is being used.

\module{8}{Model Diagnostics}
\textbf{Concepts:} residual analysis; portmanteau tests (Box--Pierce/Ljung--Box-type statistics);
overfitting; information criteria for model comparison.\\
\textbf{Learning objectives:}
\begin{itemize}
  \item Assess model adequacy using residual ACF plots and formal test statistics.
  \item Recognize signs of model misspecification and propose a correction.
  \item Compare competing candidate models using an information criterion.
\end{itemize}
\textbf{Difficulty:} \Med.

\module{9}{Forecasting}
\textbf{Concepts:} minimum mean-square-error forecasts; the general linear-process forecast formula;
prediction intervals; forecast updating as new observations arrive.\\
\textbf{Learning objectives:}
\begin{itemize}
  \item Derive point forecasts for AR, MA, ARMA, and ARIMA models.
  \item Construct prediction intervals under a normality assumption.
  \item Update a forecast when a new observation becomes available.
\end{itemize}
\textbf{Difficulty:} \MedHigh\ --- synthesizes several earlier modules; the prediction-interval
derivation is often the subtlest part.

\module{10}{Seasonal (SARIMA) Models}
\textbf{Concepts:} seasonal differencing; multiplicative SARIMA$(p,d,q)\times(P,D,Q)_s$; the full
identification/estimation/diagnostics/forecasting cycle applied to seasonal data.\\
\textbf{Learning objectives:}
\begin{itemize}
  \item Recognize seasonal structure in ACF/PACF plots and choose an appropriate seasonal
        differencing order.
  \item Specify and fit a multiplicative SARIMA model.
  \item Forecast a seasonal series and evaluate the fit via diagnostics.
\end{itemize}
\textbf{Difficulty:} \High\ --- the multiplicative notation and the interaction of seasonal and
non-seasonal orders is notationally dense.

\module{11}{Regression-Based Methods}
\textbf{Concepts:} intervention analysis; outlier detection (additive vs.\ innovational outliers);
spurious correlation; pre-whitening; cross-correlation.\\
\textbf{Learning objectives:}
\begin{itemize}
  \item Model a structural break or intervention using indicator regressors.
  \item Detect and classify outliers in a time series.
  \item Explain why serially correlated errors can produce spurious regression relationships, and how
        pre-whitening addresses this.
\end{itemize}
\textbf{Difficulty:} \MedHigh.

\module{12}{Heteroskedastic Models (ARCH/GARCH)}
\textbf{Concepts:} conditional heteroskedasticity; volatility clustering; ARCH($q$) and GARCH($p,q$)
models; basic estimation ideas.\\
\textbf{Learning objectives:}
\begin{itemize}
  \item Recognize volatility clustering as motivation for conditional-variance modeling.
  \item Write down ARCH/GARCH model equations and identify their parameters.
  \item Distinguish a conditional-\emph{mean} model (ARIMA) from a conditional-\emph{variance} model
        (ARCH/GARCH).
\end{itemize}
\textbf{Difficulty:} \High\ --- an entirely new modeling paradigm (variance modeling rather than mean
modeling), typically introduced late in the term with limited time to absorb it.

\module{13}{Additional/Advanced Topics}
\textbf{Concepts:} nonlinear time series (survey level); spectral/frequency-domain methods
(periodogram, spectral density) and their relationship to the ACF.\\
\textbf{Learning objectives:}
\begin{itemize}
  \item State the periodogram and spectral density, and relate them to the ACF (Wiener--Khinchin idea).
  \item Interpret a periodogram to detect hidden periodicities in a series.
  \item Describe, at a survey level, at least one nonlinear time series model.
\end{itemize}
\textbf{Difficulty:} \High\ --- the most mathematically abstract material in the course; covered ``as
time allows,'' so depth of coverage may vary by term.

% ================= 6. TIMELINE =================
\section{Weekly / Assessment Timeline}

\renewcommand{\arraystretch}{1.25}
\begin{longtable}{p{1.0cm} p{2.5cm} p{1.6cm} p{5.7cm} p{2.7cm}}
\toprule
\textbf{Week} & \textbf{Dates} & \textbf{Meets} & \textbf{Unit / Suggested Focus} & \textbf{Milestone} \\
\midrule
\endhead
1  & Jan.\ 19--25  & Wed.\ only (MLK holiday Mon.) & Unit I --- Module 1 begins & \\
2  & Jan.\ 26--Feb.\ 1  & Mon., Wed. & Unit I --- Module 1 cont'd / Module 2 & HW1 window opens \\
3  & Feb.\ 2--8    & Mon., Wed. & Unit I --- Module 2 cont'd / Module 3 & \\
4  & Feb.\ 9--15   & Mon., Wed. & Unit I --- Module 3 cont'd; review & \\
5  & Feb.\ 16--22  & Mon., Wed. & Unit I wrap-up & \textbf{Midterm 1} (Wed.\ Feb.\ 18) \\
6  & Feb.\ 23--Mar.\ 1 & Mon., Wed. & Unit II --- Module 4 & \\
7  & Mar.\ 2--8    & Mon., Wed. & Unit II --- Module 4 cont'd / Module 5 & \\
8  & Mar.\ 9--15   & Mon., Wed. & Unit II --- Module 5 cont'd / Module 6 & \\
   & Mar.\ 16--20  & \emph{Spring Break --- no class} & & \\
9  & Mar.\ 23--29  & Mon., Wed. & Unit II --- Module 7; review & \textbf{Midterm 2} (Wed.\ Mar.\ 25) \\
10 & Mar.\ 30--Apr.\ 5 & Mon., Wed. & Unit III --- Module 8 & \\
11 & Apr.\ 6--12   & Mon., Wed. & Unit III --- Module 9 & \\
12 & Apr.\ 13--19  & Mon., Wed. & Unit III --- Module 10 & \\
13 & Apr.\ 20--26  & Mon., Wed. & Unit III --- Module 11; review & \textbf{Midterm 3} (Wed.\ Apr.\ 22) \\
14 & Apr.\ 27      & Mon.\ only (last day of classes) & Modules 12--13, as time allows; course wrap-up & \\
\bottomrule
\end{longtable}

\textit{Homework due dates are not fixed by the syllabus (``usually due on Fridays''); HW1--HW6 are
distributed roughly every two weeks across the term, with exact dates posted on CANVAS/Gradescope. One
penalty-free late submission per semester is allowed; other late work is accepted up to 24 hours with
a 20\% penalty.}

% ================= 7. DIFFICULTY SUMMARY =================
\section{Difficulty Summary \& Preparation Tips}

\renewcommand{\arraystretch}{1.3}
\begin{longtable}{p{3.8cm} p{6.0cm} p{5.3cm}}
\toprule
\textbf{Topic} & \textbf{Why It's Difficult} & \textbf{Suggested Preparation} \\
\midrule
\endhead
Strict vs.\ weak stationarity (Module 2) &
  The two definitions coincide for Gaussian processes but not in general; students often conflate them. &
  Review joint-distribution vs.\ moment-based definitions from the probability prerequisite. \\
Backshift-operator algebra \& stationarity/invertibility roots (Module 4) &
  Requires treating an operator like a polynomial variable and reasoning about roots outside/inside the
  unit circle. &
  Refresh polynomial root-finding and geometric series; work several AR(1)/AR(2) examples by hand. \\
ACF/PACF/EACF-based identification (Module 6) &
  There is no single deterministic rule --- identification is a judgment call combining multiple
  imperfect signals. &
  Practice on many simulated series with \emph{known} true order to build pattern-recognition intuition
  before working with real data. \\
Yule--Walker / MLE estimation (Module 7) &
  Combines linear-algebra (Yule--Walker system) and calculus-based nonlinear optimization (MLE) in one
  module. &
  Review solving linear systems and basic multivariable optimization (first-order conditions) before
  this module. \\
Multiplicative SARIMA notation (Module 10) &
  Seasonal and non-seasonal orders interact multiplicatively, producing dense, easy-to-misread notation. &
  Write out the expanded backshift-operator polynomial for a simple SARIMA model by hand at least once. \\
ARCH/GARCH conditional-variance modeling (Module 12) &
  Requires a paradigm shift from modeling the conditional \emph{mean} to modeling the conditional
  \emph{variance}. &
  Before this module, review conditional expectation and conditional variance from the probability
  prerequisite. \\
Spectral/frequency-domain methods (Module 13) &
  The most abstract material in the course, often compressed into limited end-of-term time. &
  Read ahead in the relevant textbook chapter; the periodogram is easiest to grasp via a worked
  Fourier-series example. \\
\bottomrule
\end{longtable}

% ================= 8. GRADING SUMMARY =================
\section{Grading \& Policy Summary}

\begin{itemize}
  \item \textbf{Homework (25\%):} six assignments, submitted via Gradescope; problems must be tagged
        (5-point penalty for failure to tag); one penalty-free late submission per semester; otherwise
        late work accepted up to 24 hours with a 20\% penalty.
  \item \textbf{Midterms (75\%, 25\% each):} three in-class, non-cumulative exams (Feb.\ 18, Mar.\ 25,
        Apr.\ 22); formula sheet allowed (one-sided for Midterm 1, two-sided for Midterms 2--3).
  \item \textbf{Curve:} the class mean is roughly the B/B$+$ boundary; approximately $+1$SD $\to$
        A$-$/A range, $-1$SD $\to$ B$-$/C$+$ range.
  \item \textbf{Missed exams:} documented illness $\to$ grade extrapolated from the next exam; two or
        more documented missed exams $\to$ oral exam determines the course grade.
  \item \textbf{Regrades:} open for 7 days after grades are released on Gradescope (may be shorter at
        term's end).
  \item \textbf{Generative AI:} permitted for concept clarification, coding help, and brainstorming;
        submitted solutions must reflect the student's own understanding and independent work.
\end{itemize}

\vspace{1.5em}
\hrule
\vspace{0.5em}
{\footnotesize This breakdown is organized from the EN.553.439/639 syllabus (Spring 2026, Jan.\ 20,
2026 draft) and the required textbook, Cryer \& Chan, \emph{Time Series Analysis with Applications in
R}, 2nd edition. Module-to-week pacing is an illustrative distribution consistent with the syllabus's
own topic ordering and confirmed exam dates; it is \textbf{not} an official lecture-by-lecture
schedule. Weekly calendar dates are cross-checked against the published JHU Spring 2026 academic
calendar. Consult CANVAS for the authoritative lecture-by-lecture schedule and exact homework due
dates.}

\end{document}